
\chapter{Introdução}
\label{cap:intro}

O COVID-19 teve um impacto significativo em todo o mundo, resultando em hospitalizações e óbitos em grande escala. A declaração da Organização Mundial da Saúde (OMS) em março de 2020 reconheceu a gravidade da situação e a necessidade de ação global coordenada. No Brasil, o país foi amplamente afetado pela pandemia, com altos números de casos e óbitos. De acordo com a publicação "Violações dos direitos humanos no Brasil: Denúncias e Análises no Contexto da Covid-19" que pode ser encontrada no site do Conselho Nacional de Saúde pelo Ministério da Saúde, vários fatores, como a desigualdade social e regional, são responsáveis por essa quantidade de casos e óbitos. Estudos mostraram que a vacinação desempenhou um papel crucial na redução de óbitos relacionados à COVID-19, oferecendo proteção significativa, especialmente entre os idosos. As vacinas têm sido eficazes na prevenção de casos graves e na proteção da vida das pessoas, contribuindo para o controle da pandemia no Brasil.

\section{Objetivos}

Durante a execução deste trabalho, estabelecemos uma série de objetivos para aprofundar nosso entendimento sobre a relação entre vacinação e óbitos, assim como para destacar a importância da vacinação. Nossos principais objetivos são:

\begin{enumerate}
    \item[1.] Realizar um mapeamento sistemático para investigar modelos utilizados para predizer óbitos da COVID-19. 
    \item[2.] Explorar base de dados com informações de pacientes com COVID-19 na Região Centro-Oeste brasileira.
    \item[3.] Predizer casos de óbitos por COVID-19 no Estado de Goiás utilizando um modelo proposto no mapeamento sistemático.
\end{enumerate}


Ao atingir esses objetivos, esperamos contribuir para o avanço do conhecimento científico sobre a eficácia da vacinação na redução dos óbitos decorrentes da COVID-19. Nossos resultados poderão fornecer informações valiosas para os responsáveis pela saúde pública, decisores políticos e profissionais da área, ajudando a embasar estratégias efetivas no combate à pandemia.

\section{Metodologia}
% Neste trabalho foi conduzido um mapeamento sistemático para investigar o estado da arte no contexto de modelos utilizados para predizer óbitos da COVID-19. 

% 1- Mapeamento sistemático dos tipos de modelos utilizados para predição de óbitos.
% 2- Definição de modelos de predição de óbitos (com/sem a informação dos vacinados).
% 3-   
Para que fosse possível atender os objetivos do presente trabalho, as seguintes etapas foram definidas:
\begin{enumerate}
    \item[1.]Realizar um mapeamento sistemática abrangente da literatura científica existente sobre modelos de predição envolvendo vacinação e óbitos. Esse mapeamento nos permitirá analisar as abordagens utilizadas, as variáveis consideradas e as metodologias empregadas na previsão de óbitos relacionados à COVID-19 após o início da vacinação.

    \item[2.]Coletar dados da região Centro-Oeste do Brasil, específicos para cada estado, e extrair informações relevantes sobre a vacinação e os óbitos relacionados à COVID-19. Esses dados serão fundamentais para uma análise detalhada do impacto da vacinação nos índices de mortalidade em nível estadual.

    \item[3.]Desenvolver um modelo de predição de óbitos utilizando o algoritmo LSTM (Long Short-Term Memory). Esse modelo nos permitirá simular um cenário hipotético em que a população não tenha sido vacinada, possibilitando uma análise do impacto potencial da vacinação na redução dos óbitos relacionados à COVID-19. Serão considerados diferentes níveis de cobertura vacinal para avaliar seu efeito nos índices de mortalidade.
    
\end{enumerate}

\section{Organização do trabalho}
Este trabalho está organizado da seguinte forma: no Capítulo \ref{cap:background} é apresentado uma visão geral dos conceitos relacionados com o trabalho, que são essenciais para o entendimento básico do leitor, introduzindo termos científicos da área.Além disso, são apresentados gráficos que detalham o estudo do caso da Covid-19 na Região Centro-Oeste do Brasil, com informações sobre a distribuição mensal de vacinas e os tipos de dose administradas.

No Capítulo \ref{cap:mapeamento} é apresentado o protocolo do mapeamento sistemático de artigos que abordam a predição de óbitos relacionados à Covid-19 no contexto da vacinação.

A descrição do modelo de predição está presente no Capítulo \ref{cap:modelo}. São apresentadas informações sobre a base de dados utilizada no algoritmo, a seleção do modelo de predição e seus parâmetros. Além disso, os gráficos que evidenciam os resultados da avaliação desse algortimo LSTM também se encontram nesse capítulo. Por fim, as conclusões são encontradas no Capítulo \ref{cap:conclusão}.

